\begin{table}[htp]
\centering
\footnotesize
\setlength{\tabcolsep}{5pt}
\begin{tabular}{p{2.2cm} p{4.5cm} p{2cm} p{4cm}}
\toprule
\textbf{Level} & \textbf{Refinement} & \textbf{Intervention Point} & \textbf{Disturbance Coverage} \\
\midrule

L1 -- Reactive 
& Joining, departure, and role changes must be announced. Participation remains uncontrolled but becomes structurally visible. 
& During/after impact 
& Improves detection of interdependency cascades \newline 
limited influence on evolution or temporary inability \\

L2 -- Proactive 
& Joining and departure are delayed through enforced notice periods, buffering structural change and reducing instability. 
& Before impact 
& Mitigates disruption from constituent evolution and temporary inability \newline
reduces cascade severity \\

L3 -- Controlled 
& Participation becomes conditional, with enforceable obligations, admission approval, and regulated withdrawal. 
& Before formation 
& Prevents unstable integration and abrupt withdrawal \newline
strongly addresses evolution-driven disturbances \\

L4 -- Adaptive 
& Responsibilities may be reassigned to preserve capability, enabling structural reconfiguration under change. 
& Before structural conditions 
& Actively compensates for structural loss \newline
mitigates interdependencies, evolution, and partial capability degradation \\

\bottomrule
\end{tabular}
\caption{Reliability Levels}
\end{table}